\chapter{Ridimensionamento}
A volte possiamo avere la necessità di gestire strutture dati dinamiche ossia che siano in grado di adattare la loro dimensione a fronte di inserimenti e cancellazioni\\
\begin{itemize}
    \item può capitare di aver allocato una certa quantità di memoria per poi accorgersi, durante l'esecuzione del programma, che non ci basta.
    \item E' anche opportuno evitare spechi di memoria ossia, dopo aver rimosso molti elementi, ridurre la memoria allocandone una più piccola in cui memorizzare elementi rimasti
\end{itemize}
\noindent \textbf{Strategia per la gestione dinamica}\\
\begin{itemize}
    \item Alloca una nuova struttura più grande/più piccola
    \item Copia tutti gli elementi della vecchia struttura dati nella nuova
\end{itemize}
\noindent \textbf{Osservazioni}\\
\begin{itemize}
    \item La strategia appena descritta ha un costo proporzionale all'allocazione della nuova struttura + il costo della copia + il costo dell'inserimento/cancellazione
    \item Nella maggior parte delle strutture dati, il costo dominante è quello della copia
    \item Se il ridimensionamento avviene una tantum il suo costo può essere ammortizzato e distribuito sulle operazioni di inserimento successive
    \item in questo caso parleremo quindi di analisi ammortizzata del costo
    \item Consideriamo una struttura dati tabellare e vediamo
    \begin{itemize}[label=--]
        \item come sia possibile aggiungere/togliere un elemento in tempo ammortizzato costante $O(1)$, benché tali operazioni abbiano un costo unitario maggiore
        \item come sia possibile ridimensionare garantendo che la memoria inutilizzata sia sempre inferiore ad una frazione costante della memoria allocata
    \end{itemize}
\end{itemize}
\section{tabelle dinamiche}
Una tabella è una generica struttura dati che viene rappresentata in una locazione di memoria e che è divisa in celle.\\
Una tabella \textbf{T} abbia i seguenti attributi:\\
\begin{itemize}
    \item un puntatore \textbf{T.table} alla memoria allocata;
    \item \textbf{T.size} che rappresenta la dimensione della tabella
    \item un intero \textbf{T.num} che rappresenta il numero di elementi presenti nella tabella
\end{itemize}
\noindent Le operazioni che possiamo eseguire su una tabella sono:\\
\begin{itemize}
    \item \textbf{Table(T)} che crea una nuova tabella vuota T
    \item \textbf{Insert(T,x)} che inserisce l'oggetto x nella tabella T
    \item \textbf{Delete(T,x)} che rimuove l'oggetto x dalla tabella T
\end{itemize}
\noindent \subsection{Ridimensionamento: espansione}
Consideriamo dapprima il caso in cui vengano eseguite soltanto inserimenti di nuovi elementi e nessuna rimozione\\
Definiamo il fattore di carico della tabella come il rapporto $\alpha = num / size$\\
Quando $\alpha = 1$ la tabella è piena ed occorre espanderla allocando nuova memoria.\\
Un tipico metodo euristico consiste nell'allocare una nuova tabella che ha il doppio di celle della vecchia il che garantisce un fattore di carico $\alpha > 1/2$.\\ 
L'operazione di inserimento dovrà quindi essere realizzata come segue
\begin{verbatim}
INSERT(T, x)
1 If T.size == 0
2   Alloca una nuova tabella di dimensione 1
3 If T.size == T.num
4   Alloca T’ di dimensione 2*T.size e assegna al puntatore p
5   Copia tutti gli elementi di T in T’
6   Rilascia la memoria occupata da T
7   Aggiorna puntatore di T a p
8   T.table=T'
9   T.size=2*T.size
10   inserisci x in T’
11  T.num = T.num+1
\end{verbatim}
\noindent Costo:\\
Assumiamo che l'allocazione e il rilascio di memoria abbiano costo costante mentre la copia dipende dalla dimensione originale della struttura dati\\
La procedura di INSERT costerà quindi:\\
\begin{itemize}
    \item $O(1)$ se non devo espandere la tabella
    \item $O(T.size)$ se devo procedere con l'espansione
\end{itemize}
\noindent \textbf{Analisi del costo ammortizzato}\\
Consideriamo il costo di una sequenza di n \textbf{Insert} eseguite a partire da una tabella vuota\\
\begin{itemize}
    \item Il costo della k-esima Insert è 1 se non vi è espansione ed è k se vi è espansione
\end{itemize}
\noindent Siccome T.size è sempre una potenza di 2 e l'espansione si ha quando T.num=k-1 è uguale a T.size, il costo della k-esima Insert è:\medskip\\
$\begin{cases}
\text{k se k-1 è una potenza di 2} \\
\text{1 negli altri casi}
\end{cases}$\medskip\\
\noindent e il costo totale di n Insert è:\\
$$C(n)=\sum_{k=1}^n c_k=n+\sum_{j=0}^{\lfloor \log n \rfloor}2^j<n+2n=3n$$\\
\noindent se adesso dividiamo il costo totale di n operazioni per il numero di operazioni otteniamo il costo ammortizzato\\
$O(3n)/n=O(1)$
\subsection{Ridimensionamento: Contrazione}
Quando cancelliamo elementi, la struttura dati usa sempre più memoria di quanta in realtà gliene serve.\\
Per evitare uno spreco eccessivo di memoria è opportuno contrarre la tabella quando il fattore di carico $\alpha = num/size$ diventa troppo piccolo.\\
La contrazione della tabella è simile all'espansione: quando il numero di elementi in essa diventa troppo piccolo, allochiamo una nuova tabella più piccola e poi copia gli elementi della vecchia in quella nuova. Possiamo rilasciare lo spazio in memoria per la vecchia tabella restituendo al sistema di gestione della memoria.
\subsubsection{Strategia naive}
Dimezzare la memoria quando il fattore di carico $\alpha$ diventa $\le 1/2$\\
Questo ci assicura un fattore di carico $\alpha > 1/2$ anche in presenza di \textbf{Delete}\\
Purtroppo con questa strategia il costo ammortizzato delle operazioni non è più costante.\\
Consideriamo una successione costituita da $2^k$ \textbf{Insert} seguite da $2^{k-2}$ gruppi di quattro operazioni\\
\begin{center}
    \textbf{Insert - Delete - Delete - Insert}
\end{center}
per un totale di \textbf{$n=2^{k+1}$} operazioni.\\
In ogni gruppo la prima Insert comporta una espanzionedi costo $2^k$ e la seconda Delete una contrazione di costo $2^
k$.\\
Quindi ogni gruppo ha costo $\ge 2^{k+1}$\\
e in totale i $2^{k-2}$ gruppi hanno costo:
$$\ge 2^{k-2}2^{k+1}=\frac{n^2}{8}=\Omega(n^2)$$\\
e il costo ammortizzato è \textbf{$\Omega(n^2)/n=\Omega(n)$}.\\
Utilizziamocome criterio per ridimensionare $\alpha=1/2$ le sequenze di contrazioni non mi consentono di ammortizzare i costi degli inserimenti come vorrei.\\
Per le contrazioni quindi ho bisogno di svuotare a sufficienza la tabella per bilanciare il costo degli inserimenti\\
occorre quindi aspettare che $\alpha$ diventa $1/4$.
